\par\noindent\textbf{(1)} 由于(n·σ)²=I，其本征值为±1，故能级E\_+=−ℏω0/2（自旋沿n）和E\_−=+ℏω0/2。绝热演化要求磁场方向变化速率远小于能隙ω0，并且轨迹不穿越ω0=0简并点。
\par\noindent\textbf{(2)} 选择沿n的本征自旋or|n,+>=(cosθ/2,e\textasciicircum{}\{iφ\}sinθ/2)\textasciicircum{}T，则Berry联络A\_φ=i<n,+|∂\_φ|n,+>=−sin²(θ/2)=−(1−cosθ)/2。不同规范可以改变局域A\_φ，但闭路相位模2π不变。
\par\noindent\textbf{(3)} 固定θ且φ从0增至2π，Berry相位γ\_+=∮A\_φdφ=−π(1−cosθ)=−Ω\_solid/2，Ω\_solid=2π(1−cosθ)为Bloch球包围固角；相反本征支相位为+Ω\_solid/2（模2π），表现自旋1/2的球面几何性质。
\par\noindent\textbf{(4)} θ→0使Ω\_solid→0，几何相位消失；反向绕行改变闭路取向，γ变号。若Ω不远小于ω0，系统可在两个瞬时本征态间跃迁，简单Berry相位需改用完整旋转坐标精确演化而不能套绝热公式。
